\documentclass{beamer}
\usepackage{amsmath}
\usepackage{amsfonts}
\usepackage{amssymb}
\usepackage{yhmath}
\usepackage{amsthm}
\usepackage{mathrsfs}
\usepackage{mathtools}
\usepackage{mathalfa}
\usepackage{upgreek}
\usepackage{yfonts}
\usetheme{Warsaw}
%\usetheme{Madrid}
%\usetheme{Bergen}
%\usetheme{AnnArbor}
%\usetheme{CambridgeUS}\setbeamercovered{dynamic}
%\usetheme{Hannover}
\usetheme{Antibes}
%\usetheme{Copenhagen}
%\usetheme{Marburg} 
%\usefonttheme{serif}
\begin{document}
\title{Independence of the $CH$ with gentle introduction to Forcing}

\author{Rouholah Hoseini Naveh}
%{
%\setbeamercolor{background canvas}{bg=yellow}
\begin{frame}
\date{}
\maketitle

\begin{center}
1402/12/13
\end{center}
\end{frame}
\begin{frame}
\frametitle{What is $\omega$?}
$ \emptyset$ exists\\
If $A$ exists, then $\{A\}$ so\\
If $A, B$ exists, then $A \cup B$ so
\pause
\begin{itemize}
\item $\emptyset := 0$
\pause
\item $0 \cup \{0\}=\emptyset \cup \{\emptyset\}=\{\emptyset\} :=1$
\pause
\item  $1 \cup \{1\}=\{0, 1\} = \{\emptyset, \{\emptyset\}\}:=2$
\pause
\item $2 \cup \{2\}=\{0, 1, 2\}=\{\emptyset, \{\emptyset\}, \{\emptyset, \{\emptyset\}\}\}:=3$
\pause
\item $\{0, 1, 2, 3\} :=4$
\pause
\item $\{0, 1, 2, 3, \dots, n\}:=n+1$ \hspace{20 mm} $n \cup \{n\} :=n+1$
\end{itemize}
\pause
\begin{definition}
$\{0, 1, 2, 3, \dots\}:=\omega$
\end{definition}
\end{frame}
\begin{frame}
\frametitle{misunderstanding}
$\omega \cup \{\omega\}:= \omega +1$\\
\pause
$\omega+1 =\{0, 1, 2, \dots , \omega\}$
\pause
\begin{block}{}
$f: \omega +1 \longrightarrow \omega$ by $f(\omega)= 0$ and $f(n)=n+1$ is $1-1$ and onto.
\end{block}
\pause
\begin{block}{}
there is no $1-1$, onto and order preserving function
\end{block}
\pause
\vspace{10mm}
$0< 1< 2< \dots< \omega< \omega +1< \omega +2< \dots< \omega+\omega=\omega . 2 < \dots< \omega.3 < \dots< \omega . \omega = \omega^2< \dots< \omega^3< \dots< \omega^{\omega}< \dots< \omega^{\omega^{\omega}}< \omega^{\omega^{\dots}}$
\end{frame}
\begin{frame}
\begin{definition}
For each set $A$, we say $A$ is finite, if there are $n \in \omega$ and a bijection $f:n \longrightarrow A$
\end{definition}
\pause
\begin{definition}
For each infinte set $A$, we say
\begin{itemize}
\item  $A$ is countable, if there is  a bijection $f:\omega \longrightarrow A$
\item $A$ is uncountable if there is no bijection $f:\omega \longrightarrow A$
\end{itemize}
\end{definition}
\vspace{5mm}
\begin{center}
\textbf{{\Large Is there an uncountable set?}}
\end{center}
\pause
\begin{center}
\textbf{{\Large $\mathcal{P}(\omega)=\{x \colon x  \subseteq \omega\} \sim2^{\omega}= \prescript{\omega}{}{2}=\{f: \omega \longrightarrow 2\}$}}
\end{center}
\end{frame}
\begin{frame}
\frametitle{CONSISTENCY AND INDEPENDENCE}
It all started on that day in December 1873 when Georg Cantor established that \textit{\textbf{{\large the continuum is not countable}}}
\pause
\begin{theorem}
There is a bijection $f : \mathcal{P}(\omega) \longrightarrow \mathbb{R}$
\end{theorem}
\pause
\begin{theorem}
Let $A$ be a set. There is no bijection $f:A \longrightarrow \mathcal{P}(A)$
\end{theorem}
\pause
\begin{corollary}
\begin{enumerate}
\item $\mathbb{R}$ is uncountable.\\
proof: There is no bijection $f: \omega \longrightarrow \mathbb{R}$
\item There are infinitely many different uncountable.\\
$\mathcal{P}(\omega), \mathcal{P}(\mathcal{P}(\omega)), \dots $
\end{enumerate}
\end{corollary}
\end{frame}
\begin{frame}
\frametitle{The continuum hypothesis}
$\omega_1$ := the first ordinal $ \alpha$ such that there is no bijection $f:\omega \longrightarrow \alpha$
\pause
$\omega_2$ := the first ordinal $ \alpha$ such that there is no bijection $f:\omega_1 \longrightarrow \alpha$
\pause
$\omega_3$ := the first ordinal $ \alpha$ such that there is no bijection $f:\omega_2 \longrightarrow \alpha$\\
$\omega_4, \omega_5, \dots $
\pause
$$\aleph_0, \aleph_1, \aleph_3, \dots, \aleph_{\omega}, \aleph_{\omega+1}, \dots $$
\begin{block}{$c = \aleph_?$}
$f: \omega_{?} \longrightarrow \mathbb{R}$ is a bijection.
\end{block}
\end{frame}
\begin{frame}
\frametitle{Naive set theory and Axiomatic set theory}
\begin{definition}
A theory $T$ is said to be \textbf{consistent} if it doesn't proof a contradiction.\\
There is no statement $A$ such that $T \vDash A$, and $T \nvDash \neg A$.
\end{definition}
\pause
\begin{center}
{\Large  any property $P(x)$ may be used to form a set $x$}
\end{center}
\pause
\begin{center}
{\Large Let $V$ be the collection of all sets.}
\end{center}
\end{frame}
\begin{frame}
\begin{center}
{\Huge ZFC}
\end{center}
\pause
\begin{center}
{\Large Let $V$ be the collection of all sets.}
\end{center}
\pause
\begin{theorem}
\begin{enumerate}
\item if $T$ is syntactically consistent, then $T$ has a model.
\pause
\item Doesn't matter how many axiom you are using, there will always be true statements that cannot be proved.
\pause
\item If $T$ is consistent, then we can not prove this consistency in $T$.
\end{enumerate}
\end{theorem}
\end{frame}
\begin{frame}
\frametitle{consistency of $CH$ and $\neg CH$}
\begin{center}
\textbf{G\"{o}del 1939:} \ $Con(ZF) \ implies \ Con(ZFC + CH)$
\end{center}
\vspace{5mm}

\begin{center}
\textbf{Cohen 1963:} \ $Con(ZFC) \ implies \ Con(ZFC + \neg CH)$
\end{center}

\pause
\begin{center}
$$V[G] \models ZFC + c \geq \aleph_2$$
\end{center}
\end{frame}
\begin{frame}
\frametitle{FORCING}
\begin{definition}
$\langle P, \leq \rangle$ is a partial ordered set if for every $p, q, r \in P$:
\begin{enumerate}
\item $p \leq p$
\item $p \leq q$ and $q \leq p$ implies $p=q$
\item $p \leq q$ and $q \leq r$ implies $p \leq r$
\end{enumerate}
\end{definition}
\pause
\begin{block}{Forcing Notion}
Let $V \models ZFC$,\\
$\mathbb{P}= \langle P, \leq_P \rangle \in V$ a poset with largest element is a forcing notion.

$$Add(\aleph_0, 1)= \{p \colon \omega \longrightarrow 2 \ | \text{p is finite}\}$$
\end{block}
$p_1 \leq p_2$ \quad \textbf{{\large or}}\quad $p_1$ extends $p_2$ \quad \textbf{{\large or}}\quad $p_1$ is stronger than $p_2$

\textbf{{\large or}}\quad $p_1$ has more information than $p_2$ \ \textbf{{\large iff}} \ $p_2 \subseteq p_1$

\end{frame}
\begin{frame}
\vspace{-2mm}
\begin{block}{Dense subsets}
$D \subseteq \mathbb{P}$ is dense if
 $$\forall p \in \mathbb{P} \exists d \in D (d \leq p)$$
\end{block}
Fix $n \in \omega$, let $D_n=\{p \in Add(\aleph_0, 1) \colon n \in dom(p)\}$\\
\pause
let $p \in Add(\aleph_0, 1)$ arbitrary,\\
\pause
set $q= p \cup \{(n, 0)\}$,\\
\pause
clearly $q \leq p$ and $q \in D$.
\end{frame}
\begin{frame}

\begin{definition}
$G \subseteq \mathbb{P}$ is a filter if
\begin{enumerate}
\item For all $p, q \in G$, there is $r \in G$ such that $r \leq p$ and $r \leq q$
\item  If $p \in G$, then all $q \in \mathbb{P}$ with $p \leq q$ are in $G$.
\end{enumerate}
\end{definition}
\pause
\begin{block}{Generic Set}
$G$ is Generic if
\begin{itemize}
	\item $G \subseteq \mathbb{P}$ is a filter
	\item $G$ meets every dense subset of $\mathbb{P}$ that lies in $V$
\end{itemize}
\end{block}
\end{frame}
\begin{frame}
\frametitle{Magics}
\begin{theorem}
If $G \subset Add(\aleph_0, 1)$ is generic, then $G \notin V$ 
\end{theorem}
\pause
\begin{block}{spotlight}
Note that $G \subset V$\\
Let $V$ be a (nice)model and $ON$ be the class of all ordinals.\\
$V \cap ON$ is an ordinal but $V \cap ON \notin V$.
\end{block}
\pause
We can adjoin $G$ to $V$ by an amazing method and get $V[G]$
\begin{enumerate}
\item $V[G] \models ZFC$
\item $V \subseteq V[G]$ and $G\in V[G]$
\item $V[G] \cap ON=V \cap ON$
\item If $M$ satisfies $(1)-(3)$, then $V[G] \subseteq M$.
\end{enumerate}
\end{frame}
\begin{frame}
\frametitle{more magics}
\begin{theorem}
For every formula $\varphi(x_1, \dots, x_n)$ we can find a formula $\varphi'(x_1, \dots, x_n, p, \mathbb{P})$ such that\\
for every forcing notion $\mathbb{P}$ and every $p \in \mathbb{P}$\\
$V \models \varphi'(x_1, \dots, x_n,p, \mathbb{P})$ iff for every generic $G \subseteq \mathbb{P}$ with $p \in G$\\
$V[G] \models \varphi(x_1, \dots, x_n)$
\end{theorem}
\end{frame}
\begin{frame}
Working in $V[G]$:
\begin{theorem}
$\bigcup G : \omega \longrightarrow 2$ is a function in $V[G]$
\end{theorem}
\pause
\begin{theorem}
For every $f:\omega \longrightarrow 2$ in $V$, $f \neq \bigcup G$
\end{theorem}
\pause
\begin{proof}
Let $f$ as above,\\
\pause
consider the dense set $D_f=\{p \in Add(\aleph_0, 1) \colon \exists n \in dom(p) (p(n) \neq f(n))\}$\\
\pause
Since $G$ is generic, so $G \cap D_f \neq \emptyset$\\
\pause
hence there exists $n$ such that $\bigcup G (n) \neq f(n)$
\end{proof}
\end{frame}
\begin{frame}
\frametitle{Cohen forcing}
$$Add(\aleph_0, \aleph_2)=\{p:(\omega \times \omega_2) \longrightarrow 2 | \text{p is finite}\}$$
this forcing adds $\aleph_2$ many real.\\
\begin{block}{}
Cohen showed that $c^{V[G]} \geq \aleph_2^{V}$
\end{block}
\begin{quotation}
Does $\aleph_2^V=\aleph_2^{V[G]}$?
\end{quotation}
\end{frame}
\begin{frame}
\begin{center}
\textbf{{\Huge Thank You}}
\end{center}
\end{frame}
\end{document}
